% paper/sections/appendix_numerics.tex
%
% 付録（数値解析・安定性・圧力関連の補足証明）
%   sec:dtau_derive        — 擬似時間反復の収束率解析と最適 Δτ 導出
%   app:checkerboard_mode  — チェッカーボードモードの代数的証明
%   app:capillary_cfl      — 毛管波制約の導出
%   app:fvm_face_coeff     — PPE 面係数 a_f の FVM 積分論理
%   app:grid_ale           — 時間変化格子と ALE 効果（本稿では未実装）
%   app:godunov_ls         — 標準 LS 法の Godunov 型スキーム（参考）
%

% ═══════════════════════════════════════════════════════════════════════════════
\section{スウィープ型擬似時間反復の収束率解析と最適 \texorpdfstring{$\Delta\tau$}{Δτ} の導出}
\label{sec:dtau_derive}
% ═══════════════════════════════════════════════════════════════════════════════

本付録は §\ref{sec:ppe_pseudotime} で使用した収束率指標
$\gamma(t) = (1+t^2)/(1+t)^2$ を導出し，最適擬似時間刻み $\Delta\tau_{\mathrm{opt}}$ を
理論的に確定する．

\subsection*{設定}

スウィープ型擬似時間 PPE 反復の1ステップを
\begin{equation}
  (\mathbf{I} + \Delta\tau \mathcal{L}_x)(\mathbf{I} + \Delta\tau \mathcal{L}_y)\,
  p^{m+1} = p^m + \Delta\tau\,q_h
  \label{eq:dtau_sweep}
\end{equation}
と書く．ここで $\mathcal{L}_h = \mathcal{L}_x + \mathcal{L}_y$ は
CCD 離散 Laplacian の $x$・$y$ 方向分割であり，
それぞれ固有値 $\lambda_x,\lambda_y \ge 0$ を持つ（正定値）．

\textbf{スプリッティング不動点．}
式 \eqref{eq:dtau_sweep} の不動点 $\tilde{p}$（$p^{m+1}=p^m=\tilde{p}$）は
\[
  \bigl(\mathcal{L}_h + \Delta\tau\,\mathcal{L}_x\mathcal{L}_y\bigr)\tilde{p} = q_h
\]
を満たす．これは元の PPE $\mathcal{L}_h p^* = q_h$ とは $\mathcal{O}(\Delta\tau^2)$ だけ異なり，
この差がスプリッティング誤差（§\ref{warn:ppe_splitting}）の根源である．

\subsection*{誤差収縮の解析}

誤差 $e^m \equiv p^m - \tilde{p}$ を定義する．
式 \eqref{eq:dtau_sweep} から不動点方程式を辺々引くと
\[
  (\mathbf{I}+\Delta\tau\mathcal{L}_x)(\mathbf{I}+\Delta\tau\mathcal{L}_y)\,e^{m+1} = e^m
\]
となり，誤差は\textbf{斉次}の縮小方程式に従う．

\subsection*{スカラー化}

等方格子 $(\Delta x = \Delta y = h)$ かつ一様係数を仮定すると
$\lambda_x = \lambda_y = \lambda/2$（$\lambda$ は $\mathcal{L}_h$ の固有値）．
$t \equiv \Delta\tau\lambda/2$ とおけば，スペクトル解析から

\begin{equation}
  e^{m+1} = \frac{e^m}{(1+t)^2}
  \label{eq:error_contraction}
\end{equation}

が得られる．すなわちスウィープ型反復は，スプリッティング不動点 $\tilde{p}$ への誤差を
各反復ごとに $1/(1+t)^2$ 倍に縮小する（\textbf{無条件安定}；任意の $t > 0$ で縮小率 $< 1$）．

\medskip
\textbf{実用的な収束率指標 $\gamma(t)$ の導出．}
上記の純粋な誤差縮小率 $1/(1+t)^2$ は $t \to \infty$ で $0$ に近づくため，
「大きな $\Delta\tau$ が常に良い」と誤解される恐れがある．
しかし大きな $\Delta\tau$ はスプリッティング不動点誤差 $\|\tilde{p} - p^*\| \sim \mathcal{O}(\Delta\tau^2)$
をも拡大する．このトレードオフを捉える指標として，
等方格子・一様係数・スカラー固有値 $\lambda$（$t = \Delta\tau\lambda/2$）の設定で
$p^0 = 0$ から1スウィープ後の真の残差比を厳密に計算する．

\textbf{1スウィープ後の解：} 式 \eqref{eq:dtau_sweep} に $p^0 = 0$ を代入すると
$(1+t)^2 p^1 = \Delta\tau\,q_h$，すなわち $p^1 = \Delta\tau q_h/(1+t)^2$．
真の解は $\lambda p^* = q_h$，$\Delta\tau = 2t/\lambda$ より
$p^* = q_h/\lambda = \Delta\tau q_h/(2t)$．

\textbf{残差比の計算：}
\begin{align*}
  r^0 &= \lambda(p^0 - p^*) = -q_h \\
  r^1 &= \lambda(p^1 - p^*) = \lambda\!\left[\frac{\Delta\tau q_h}{(1+t)^2} - \frac{\Delta\tau q_h}{2t}\right]
       = \lambda\,\Delta\tau\,q_h \cdot \frac{2t - (1+t)^2}{2t(1+t)^2}
       = -q_h\cdot\frac{1+t^2}{(1+t)^2}
\end{align*}
ここで $2t-(1+t)^2 = -(1+t^2)$，$\lambda\,\Delta\tau = 2t$ を用いた．したがって：

\begin{equation}
  \gamma(t) \equiv \frac{|r^1|}{|r^0|} = \frac{1 + t^2}{(1+t)^2}
  \label{eq:gamma_t}
\end{equation}

（$t > 0$ に対して $\gamma(t) < 1$；初期残差を1スウィープで $\gamma$ 倍に縮小する指標．）
この指標は「収束速度」と「スプリッティング誤差下限」の両方を反映しており，
$t \to 0$ でも $t \to \infty$ でも $\gamma \to 1$（収束が遅い）という正しい定性的挙動を示す．

\subsection*{最適点の導出}

$\gamma(t)$ を最小化する $t^*$ を求める：
\begin{equation}
  \frac{d\gamma}{dt} = \frac{2(t-1)}{(1+t)^3} = 0
  \quad\Longrightarrow\quad t^* = 1
\end{equation}

このとき $\gamma(1) = 2/4 = 1/2$．すなわち最適擬似時間刻みでは
$\gamma$ の意味で\textbf{1スウィープごとに収束率指標が最小値}をとり，
収束と分割誤差のトレードオフが最も良くバランスする．

$t^* = 1$ を $\Delta\tau$ に戻すと $\Delta\tau_{\mathrm{opt}} = 2/\lambda_{\max}$．
CCD Laplacian の最大固有値は $\lambda_{\max} \approx 3.43\,a_{\max}/h^2$（$a_{\max}$ は最大逆密度係数）であるから

\begin{equation}
  \Delta\tau_{\mathrm{opt}} \approx \frac{2h^2}{3.43\,a_{\max}}
                             \approx \frac{0.58\,h^2}{a_{\max}}
  \label{eq:dtau_opt_derive}
\end{equation}

これは §\ref{sec:ppe_pseudotime} の経験的指針 $C \approx 1$ の理論的根拠となる．

\subsection*{$t \to 0$ および $t \to \infty$ の挙動}

\begin{itemize}
  \item $\gamma(t) \to 1$（$t \to 0$）：$\Delta\tau$ が小さすぎると誤差縮小率 $1/(1+t)^2 \to 1$ となり，
        スプリッティング不動点 $\tilde{p}$ への収束が遅い．
  \item $\gamma(t) \to 1$（$t \to \infty$）：$\Delta\tau$ が大きすぎるとスプリッティング不動点誤差
        $\|\tilde{p} - p^*\| \sim \mathcal{O}(\Delta\tau^2)$ が拡大し，真の解 $p^*$ からの乖離が増す．
  \item 最適点 $t=1$（$\Delta\tau = \Delta\tau_{\mathrm{opt}}$）：指標 $\gamma$ が最小値 $1/2$ をとり，
        収束速度と分割誤差のトレードオフが最良にバランスする．
        式~\eqref{eq:error_contraction} の純粋な誤差縮小率はこの点で $1/4$ であり，
        収束は速いが分割誤差下限が残る（残差収束判定により実用上は問題ない；§\ref{warn:ppe_splitting} 参照）．
\end{itemize}
\qed

% ═══════════════════════════════════════════════════════════════════════════════
\section{チェッカーボードモードの代数的証明}
\label{app:checkerboard_mode}
% ═══════════════════════════════════════════════════════════════════════════════

格子間隔 $h$ の一様格子上でチェッカーボードモード
\begin{equation}
  p_j = P_0(-1)^j
\end{equation}
を考える．中心差分による圧力勾配は：
\[
  \left.\frac{dp}{dx}\right|_j \approx \frac{p_{j+1} - p_{j-1}}{2h}
  = \frac{P_0(-1)^{j+1} - P_0(-1)^{j-1}}{2h}
  = \frac{-P_0(-1)^j - (-P_0(-1)^j)}{2h} = 0
\]
すなわち，$p_j = P_0(-1)^j$ は（定数圧力を除いて）あらゆる物理的な圧力から区別できない非物理的なモードであり，
中心差分では\textbf{零勾配を与える（ゴーストモード）}となる．

このモードが圧力場に混入すると，速度場を修正しないまま圧力のみ発散する「チェッカーボード不安定」を引き起こす
（対策は §\ref{sec:rhie_chow} および §\ref{sec:balanced_force} で詳述している）．

% ═══════════════════════════════════════════════════════════════════════════════
\section{毛管波制約の導出：分散関係から CFL 条件へ}
\label{app:capillary_cfl}
% ═══════════════════════════════════════════════════════════════════════════════

\textbf{Step 1：毛管波の分散関係}

気液界面上の微小振幅波（毛管波）の線形安定解析から，
角周波数 $\omega$ と波数 $k$ の関係（分散関係）は：
\begin{equation}
  \omega^2 = \frac{\sigma\,k^3}{\rho_l + \rho_g}
  \label{eq:capillary_dispersion_app}
\end{equation}
位相速度 $c_\sigma = \omega/k = \sqrt{\sigma k/(\rho_l+\rho_g)}$，
群速度 $c_g = d\omega/dk = (3/2)c_\sigma$．

\textbf{Step 2：格子上の最大分解可能波数}

格子間隔 $\Delta x$ の離散格子で分解できる最高波数はナイキスト波数
$k_{\max} = \pi/\Delta x$ である．

\textbf{Step 3：CFL 条件の適用}

$k_{\max}$ に対応する群速度で CFL 条件 $c_g(k_{\max})\,\Delta t \le \Delta x$ を適用する：
\begin{align*}
  c_g(k_{\max}) &= \frac{3}{2}\sqrt{\frac{\sigma\pi/\Delta x}{\rho_l+\rho_g}} \\
  \Delta t &\le \frac{\Delta x}{c_g(k_{\max})}
  = \frac{2}{3}\sqrt{\frac{(\rho_l+\rho_g)\,\Delta x^3}{\sigma\pi}}
\end{align*}
係数 $2/(3\sqrt{\pi}) \approx 0.376$ を保守的に $1/\sqrt{2\pi}$ に丸めると式 \eqref{eq:dt_sigma} を得る：
\[
  \Delta t_\sigma = \sqrt{\frac{(\rho_l+\rho_g)\,\Delta x^3}{2\pi\sigma}}
\]

\textbf{物理的解釈：} $\Delta t_\sigma \propto \Delta x^{3/2}$ は分散関係 $\omega \propto k^{3/2}$ に由来する．
格子を半分にすると $\Delta t_\sigma$ は $1/2^{3/2} \approx 0.354$ 倍（$\approx 2.8$ 倍制約強化）．

\medskip
\textbf{大密度比（$\rho_l \gg \rho_g$）での解釈：}
$\rho_l/\rho_g = 1000$（水/空気）では $\rho_l + \rho_g \approx \rho_l$（気相寄与は $0.1\%$ 以下）．
表面張力が\textbf{復元力}（ばね），液相 $\rho_l$ が\textbf{慣性}（質量）を担い，
気相の慣性は無視できる．密度比 1000 の系では密度比 1 の系より
$\sqrt{1000} \approx 32$ 倍大きな $\Delta t_\sigma$ が許される（制約緩和）．\qed

% ═══════════════════════════════════════════════════════════════════════════════
\section{\texorpdfstring{PPE 面係数 $a_f = 2/(\rho_L+\rho_R)$ の FVM 積分論理}{PPE 面係数 a\_f = 2/(ρ\_L+ρ\_R) の FVM 積分論理}}
\label{app:fvm_face_coeff}
% ═══════════════════════════════════════════════════════════════════════════════

PPE 演算子の係数 $a(x) = 1/\rho(x)$ について，
セル $P$（位置 $x_L$）と $E$（位置 $x_R$）の間の等価面係数を FVM 積分で求める．

フラックス $F = a_f(p_E - p_P)/\Delta x$ において，
連続体のフラックス $F = a(x)\,dp/dx$ を区間 $[x_L, x_R]$ で積分して面係数 $a_f$ を求める．
$a(x)$ が区間内で区分的定数（各セル内で $a_L = 1/\rho_L$，$a_R = 1/\rho_R$）のとき：
\[
  \int_{x_L}^{x_R} \frac{dx}{a(x)}
  = \frac{\Delta x/2}{a_L} + \frac{\Delta x/2}{a_R}
  \quad\Rightarrow\quad
  a_f = \frac{\Delta x}{\dfrac{\Delta x/2}{a_L}+\dfrac{\Delta x/2}{a_R}}
  = \frac{2a_L a_R}{a_L+a_R}
  = \frac{2}{\rho_L+\rho_R}
\]
これは直列抵抗の等価抵抗（$1/R_\mathrm{eq} = 1/R_1+1/R_2$）と同じ論理（調和平均）であり，
熱伝導率の界面補間で使われる標準公式と同一構造を持つ．\qed

% ═══════════════════════════════════════════════════════════════════════════════
\section{時間変化格子と ALE 効果（本稿では未実装）}
\label{app:grid_ale}
% ═══════════════════════════════════════════════════════════════════════════════

本付録は §\ref{sec:grid_gen} の格子が時間発展する場合の ALE（Arbitrary Lagrangian-Eulerian）補正を説明する．
\textbf{本稿の実装では格子を時間固定として扱う（$\bm{v}_\text{mesh} = 0$）}か，タイムステップごとに格子を再生成するが $\bm{v}_\text{mesh}$ の補正は省略する．

\begin{tcolorbox}[warnbox, title={ALE 定式化の概要と省略の条件}]
ALE 定式化では，格子速度 $\bm{v}_\text{mesh}$ を輸送方程式に追加する必要がある：
\[
  \frac{\partial\psi}{\partial t} + \nabla\cdot\bigl((\bm{u} - \bm{v}_\text{mesh})\psi\bigr) = 0
\]
ただし，この追加項の離散化は実装の複雑さを大きく増加させる．
\end{tcolorbox}

\textbf{誤差の発生機構：}
1ステップあたりの界面移動量を $\delta = |\bm{u}|\Delta t$ とすると，
ALE 補正を省略することで $\bm{v}_\text{mesh}$ による追加対流量が無視される：
\[
  \varepsilon_{\text{ALE}} \sim |\bm{v}_\text{mesh}| \cdot |\bnabla\psi| \cdot \Delta t \sim U_\Gamma \cdot \frac{1}{\varepsilon} \cdot \Delta t
\]
ここで $|\bnabla\psi| \sim 1/(4\varepsilon)$（界面近傍プロファイル），$U_\Gamma$ は界面速度である．

\textbf{近似が有効な条件：}
\begin{center}
\renewcommand{\arraystretch}{1.3}
\begin{tabular}{>{\raggedright\arraybackslash}p{48mm}c>{\raggedright\arraybackslash}p{42mm}}
\hline
条件 & 移動量 $\delta/h$ & ALE 省略の妥当性 \\
\hline
緩やかな界面変形 & $< 0.1$ & 良好（誤差 $< 10\%$ of $\Delta t^2$ 項） \\
標準的な問題 & $\approx 0.5$ & 条件付きで可（格子収束確認要） \\
大密度比 ($\rho_l/\rho_g{=}1000$) の高速界面 & $> 1$ & 精度劣化の可能性大（ALE 実装推奨） \\
\hline
\end{tabular}
\end{center}

$\rho_l/\rho_g = 1000$ の場合，液相変形が速く $\delta/h > 1$ となることがあり，ALE 補正なしの格子更新は実質的に $O(\Delta t)$ 精度（1次）となる．
格子収束テスト（第\ref{sec:verification}章）に加えて，時間刻み $\Delta t$ を半分にしたとき誤差が $1/4$ 倍（2次）か $1/8$ 倍（3次）になるかを確認し，ALE 省略の影響を定量的に評価すること．

% ═══════════════════════════════════════════════════════════════════════════════
\section{Rhie--Chow 補正と CCD 高次精度の整合性}
\label{app:rc_precision}
% ═══════════════════════════════════════════════════════════════════════════════

式~\eqref{eq:rc-face} の補正項は $O(h^2)$ の誤差を持つが，CCD の $O(h^6)$ 精度を律速しない理由を示す．

フェイス $f$（$x_P$ と $x_E$ の中点，$x_f = (x_P+x_E)/2$，$x_E - x_P = h$）での各勾配評価の
テイラー展開：

\textbf{(1) 直接フェイス差分：}
$x_P = x_f - h/2$，$x_E = x_f + h/2$ として
\[
  p_E - p_P = h\,p'(x_f) + \frac{h^3}{24}p'''(x_f) + \Ord{h^5}
\]
を $h$ で割ると
\begin{equation}
  (\nabla p)_f \equiv \frac{p_E - p_P}{h} = p'(x_f) + \frac{h^2}{24}p'''(x_f) + \Ord{h^4}
  \label{eq:rc_face_taylor}
\end{equation}

\textbf{(2) CCD 勾配の算術平均：}
CCD は各セル中心で $O(h^6)$ 精度の1階微分を返す：
$(D^{(1)}_\mathrm{CCD}\,p)_P = p'(x_P) + \Ord{h^6}$，
$(D^{(1)}_\mathrm{CCD}\,p)_E = p'(x_E) + \Ord{h^6}$．
これらの算術平均をテイラー展開：
\[
  p'(x_P) = p'(x_f) - \tfrac{h}{2}p''(x_f) + \tfrac{h^2}{8}p'''(x_f) + \Ord{h^3},
  \quad
  p'(x_E) = p'(x_f) + \tfrac{h}{2}p''(x_f) + \tfrac{h^2}{8}p'''(x_f) + \Ord{h^3}
\]
\begin{equation}
  \overline{(\nabla p)}_f \equiv \frac{1}{2}\bigl[(D^{(1)}_\mathrm{CCD}\,p)_P + (D^{(1)}_\mathrm{CCD}\,p)_E\bigr]
    = p'(x_f) + \frac{h^2}{8}p'''(x_f) + \Ord{h^4}
  \label{eq:rc_avg_taylor}
\end{equation}
（奇数次項が相殺し，$O(h^2)$ の次数で $p'(x_f)$ を近似することに注意；$O(h^6)$ ではない）

\textbf{(3) Rhie-Chow 検出項（差）：}
\begin{equation}
  (\nabla p)_f - \overline{(\nabla p)}_f
  = \left(\frac{1}{24} - \frac{1}{8}\right)h^2 p'''(x_f) + \Ord{h^4}
  = -\frac{h^2}{12}p'''(x_f) + \Ord{h^4} = \Ord{h^2}
  \label{eq:rc_detection_taylor}
\end{equation}
補正量は $\Ord{\Delta t \cdot h^2}$．

\begin{itemize}
  \item \textbf{滑らかな圧力場：} 補正が $\Ord{h^2}$ であり，CSF モデル誤差 $\Ord{\varepsilon^2} \approx \Ord{h^2}$ と同程度．CCD の $\Ord{h^6}$ を律速する新たな誤差源とはならない．
  \item \textbf{チェッカーボード場（波長 $2h$）：} $p_E - p_P = \Ord{1}$，$\overline{(\nabla p)}_f \approx 0$ のため補正量が $\Ord{1}$ の大きな値となり非物理振動を除去する．
  \item \textbf{スペクトル選択的フィルタ：} 低波数（物理情報）は CCD $\Ord{h^6}$ で処理し，高波数（チェッカーボード）のみ $\Ord{h^2}$ で減衰する選択的フィルタとして機能する．
\end{itemize}

Balanced-Force 条件（§\ref{sec:balanced_force}）への影響はない：Rhie-Chow 補正はフェイス速度の安定化補助項であり，主方程式の CCD 演算子とは独立の役割を担う．

% ═══════════════════════════════════════════════════════════════════════════════
\section{標準 Level Set 法の Godunov 型スキーム（参考）}
\label{app:godunov_ls}
% ═══════════════════════════════════════════════════════════════════════════════

本付録は §\ref{sec:weno5} で用いる CLS 法の WENO5 スキームとは別に，
\textbf{標準レベルセット（符号付き距離関数）}移流に対して使われる Godunov 型 upwinding スキームを説明する．
本稿の主実装（CLS + WENO5）では使用しないが，実装の選択肢として参考に記す．

\subsection{Godunov 型スキームの基本形}

符号付き距離関数 $\phi$ の移流方程式 $\phi_t + \bm{u}\cdot\nabla\phi = 0$ に対して，
速度場 $\bm{u}$ の発散フリー条件 $\nabla\cdot\bm{u}=0$ のもとで保存形 $\phi_t + \nabla\cdot(\bm{u}\phi) = 0$ と等価となる．

$x$ 方向フラックスを Godunov 法で評価する：
\begin{equation}
  F_{i+1/2} = u_{i+1/2}^+ \phi_i + u_{i+1/2}^- \phi_{i+1}
\end{equation}
ここで $u^+ = \max(u, 0)$，$u^- = \min(u, 0)$ は速度の上流風分解である．

\subsection{再初期化との組み合わせ}

標準 LS 法では移流後に符号付き距離条件 $|\nabla\phi| = 1$ が崩れるため，
再初期化方程式
\[
  \phi_\tau + S(\phi^0)(|\nabla\phi|-1) = 0
\]
を仮想時間 $\tau$ 方向に数回反復する（$S(\phi^0)$ は初期符号に合わせたスムーズ符号関数）．

CLS 法（本稿主方式）では再初期化が $\psi$ の compact プロファイルを保つよう設計されており，
Godunov 型スキームと組み合わせるよりも質量保存性に優れる（§\ref{sec:reinit}）．
